% ---------------- Dedication ----------------
\chapter*{Dedication}
We dedicate this work to our parents, whose unwavering support, sacrifices, and prayers have
been the foundation of our success. Their encouragement, patience, and belief in our abilities
have guided us throughout this journey.

We also dedicate this to our teachers and mentors who inspired us to pursue knowledge with
dedication and integrity.

% ---------------- Acknowledgement ----------------
\chapter*{Acknowledgement}
All praise and gratitude are due for the successful completion of this project. The authors
would like to express their sincere appreciation to their supervisor for continuous guidance, valuable
suggestions, and constructive feedback throughout the development of this thesis. The supervision
and mentorship provided significant support in completing this work successfully.

Special thanks are extended to the faculty members of the Department of Computer Science
for their academic support and encouragement during the degree program.

The authors are also deeply grateful to their parents and families for their unwavering support,
motivation, and prayers, which have been a constant source of strength.

Finally, appreciation is extended to everyone who directly or indirectly contributed to the
successful completion of this project.

% ---------------- Executive summary ----------------
\chapter*{Executive Summary}
\addcontentsline{toc}{chapter}{Executive Summary}
Organizations keep their contracts, non-disclosure agreements, leases, employment letters and
affidavits as a mix of digital PDFs, scanned pages and handwritten documents, next to the laws and
policies those documents depend on. Simple questions such as \emph{``When was the NDA with Acme
signed?''} or \emph{``Which agreements expire this year?''} usually mean opening files one by one,
and scanned or handwritten pages cannot be searched at all. General ``chat with your PDF'' tools do not
solve this: they treat every file as loose text, keep no structured record of dates and parties,
do not control which employee may see which document, and keep no history.

\sysname{} is a web-based legal records intelligence system for an organization. Employees
upload documents into one shared, access-controlled repository. A background pipeline reads every
page (digital text, Tesseract OCR for printed scans, and a vision-language model for handwriting),
decides whether the document is a \emph{record} (a contract between parties) or \emph{reference}
material (a law book or policy), splits it along its own section structure, and indexes it for
both keyword and semantic search. For records, a language model extracts eleven structured fields
(parties, signing, effective and expiry dates, notice period, governing law, value and others); every
value carries the page and the exact quote that proves it, and values that cannot be verified go to a
human review queue.

A query router sends each question down one of three paths: \emph{lookup} questions are answered
directly from the structured records table, \emph{knowledge} questions (for example about the
Contract Act 1872) are answered from retrieved passages with numbered citations, and
\emph{search} questions return the matching documents. Every answer opens the source document at
the cited page with the evidence highlighted. Role-based access control (Admin, Manager, Employee,
Auditor) together with department and named-person sharing is enforced \emph{before} any
retrieval, so the assistant can never quote a document the user may not open. Four kinds of history
are kept: query history, document versions with comparison, a field-edit history, and an
append-only audit log. Reminders warn the right people before agreements expire or notice
deadlines pass.

A clause classifier was fine-tuned on the LEDGAR dataset (100 clause types) and the system was
evaluated on the public CUAD contract dataset. Section-aware chunking with hybrid retrieval raised
Hit@5 from 82.8\% (the previous fixed-size, semantic-only design) to 93.0\%; field extraction reached a
micro-F1 of 75.9\% with 99.4\% of evidence quotes found word-for-word on the cited page. All 26
automated tests, which include the full role and access-level matrix, pass.

The system is built with Vue.js 3, Django 6 with Django REST Framework, and PostgreSQL 16 with the
pgvector extension, and can be deployed with Docker Compose.

% ---------------- Abbreviations ----------------
\chapter*{Abbreviations}
\addcontentsline{toc}{chapter}{Abbreviations}
\begin{longtable}{p{3.2cm}p{11.5cm}}
\toprule
\textbf{Abbreviation} & \textbf{Full Form} \\ \midrule
\endhead
API & Application Programming Interface \\
BM25 & Best Matching 25 (keyword ranking function) \\
CER / WER & Character Error Rate / Word Error Rate \\
CUAD & Contract Understanding Atticus Dataset \\
DFD & Data Flow Diagram \\
DRF & Django REST Framework \\
ERD & Entity Relationship Diagram \\
FR / NFR & Functional Requirement / Non-Functional Requirement \\
FTS & Full-Text Search \\
GIN & Generalized Inverted Index \\
HNSW & Hierarchical Navigable Small World \\
JWT & JSON Web Token \\
LEDGAR & Labeled EDGAR contract-provision dataset \\
LLM & Large Language Model \\
MRR & Mean Reciprocal Rank \\
NDA & Non-Disclosure Agreement \\
OCR & Optical Character Recognition \\
PK / FK & Primary Key / Foreign Key \\
RAG & Retrieval-Augmented Generation \\
RBAC & Role-Based Access Control \\
REST & Representational State Transfer \\
RRF & Reciprocal Rank Fusion \\
SDD & Software Design Document \\
SPA & Single-Page Application \\
SRS & Software Requirements Specification \\
TF-IDF & Term Frequency--Inverse Document Frequency \\
TPM / TPD & Tokens Per Minute / Tokens Per Day (LLM API quotas) \\
UC & Use Case \\
UI & User Interface \\
UML & Unified Modeling Language \\
VLM & Vision-Language Model \\
3NF & Third Normal Form \\
\bottomrule
\end{longtable}
